\chapter{Introduction and Problem Description}
\label{chap:introduction}

\section{Problem Context}
\label{sec:context}

Managing large-scale Capital Expenditure (CAPEX) portfolios requires the coordinated execution of numerous concurrent investment projects — infrastructure expansions, technology rollouts, engineering programs — each competing for a shared, finite pool of specialized human resources. A central challenge in this environment is the lack of \textit{forward-looking transparency} about which competencies are required, in what quantity, and at which point in the project lifecycle.

Without such foresight, organizations struggle to ensure that projects are staffed with the right people at the right time. Bottlenecks emerge when demand for a specific skill — civil engineering, systems integration, project controls — peaks across multiple projects simultaneously. By the time these conflicts surface in conventional planning tools, the window for effective corrective action has already closed, leaving only costly remedies: emergency contracting, project scope reduction, or schedule slippage.

\section{Problem Statement}
\label{sec:problem}

Current resource planning in CAPEX portfolio environments is predominantly manual and reactive. Spreadsheet-based tools and generic Enterprise Resource Planning (ERP) modules lack the analytical capability to forecast skill-specific demand across a heterogeneous portfolio. As a result:

\begin{itemize}
  \item Resource bottlenecks are identified too late for proactive intervention.
  \item Allocation decisions rely on intuition rather than data-driven optimization.
  \item Strategic workforce options — hiring, contracting, upskilling — are not evaluated systematically.
\end{itemize}

This creates a structural gap between the \textit{strategic intent} of portfolio planning and the \textit{operational reality} of resource availability.

\section{Relevance and Motivation}
\label{sec:motivation}

The increasing complexity of CAPEX portfolios, combined with growing talent scarcity in technical disciplines, makes this problem both timely and consequential. Organizations that can anticipate resource constraints early gain a measurable competitive advantage: they can decide proactively whether to \textit{make} (upskill existing employees), \textit{buy} (recruit), or \textit{borrow} (contract) the required competencies — rather than reacting under pressure.

Advances in AI — particularly the convergence of \textit{symbolic AI} (knowledge representations, ontologies) and \textit{subsymbolic AI} (data-driven machine learning, generative AI) into \textit{hybrid AI} approaches — offer a promising foundation for addressing this gap. Gartner identifies hybrid AI as an emerging technology trend with significant enterprise applicability \citep{gartner2024hype}. At the same time, structured knowledge sources such as the European Skills, Competences, Qualifications and Occupations ontology (\textbf{ESCO}, \url{https://esco.ec.europa.eu}) provide a standardized, machine-readable taxonomy of skills and occupations that can serve as a semantic backbone for resource planning systems.

\section{Methodological Foundation: The AI-4-SME Framework}
\label{sec:ai4sme}

This thesis adopts the \textbf{AI-4-SME Framework} \citep{peter2025ai4sme} — developed at FHNW School of Business — as its overarching methodological structure. The framework provides a three-phase approach to identifying, building, and deploying AI solutions in organizational contexts:

\begin{description}
  \item[Phase 1 — Design] AI-relevant opportunities are identified through analysis at the company level (products and services), the process level (knowledge- and data-intensive tasks), and the task level (solution ideas via design thinking). A central instrument is the distinction between \textit{Knowledge-Intensive Tasks} (KIAs) — requiring problem solving, analysis, and decision making — and \textit{Data-Intensive Tasks} (DIAs) — requiring data analysis and transaction processing. Resource demand forecasting and bottleneck detection map directly onto KIAs, while allocation optimization maps onto DIAs.

  \item[Phase 2 — Build] The framework defines two implementation paths — a \textit{data-based path} (machine learning with frameworks such as TensorFlow or PyTorch) and a \textit{knowledge-based path} (knowledge graph systems such as Neo4j, GraphDB, Stardog, or RDFox) — which can be combined into a \textit{hybrid AI solution}. This thesis follows the hybrid path: ESCO-grounded knowledge graphs provide structured domain knowledge, while machine learning models capture temporal demand patterns from project data.

  \item[Phase 3 — Run] The deployed system is integrated into the organization's productive environment, addressing corporate governance (GDPR, EU AI Act compliance), IT infrastructure, application-level API integration with ERP/PPM systems, and change management for human-AI collaboration.
\end{description}

This thesis investigates how hybrid AI — combining ontology-based knowledge representation with data-driven forecasting and optimization, structured according to the AI-4-SME framework — can power an intelligent recommendation engine for CAPEX resource management.
